% Shared preamble for the tutorial and the technical companion.
% Keep notation aligned with guidance/NOTATION.md.
\usepackage[T1]{fontenc}
\usepackage[utf8]{inputenc}
\usepackage{lmodern}
\usepackage[scale=.9]{tgheros}
\usepackage[margin=.85in,headheight=14pt]{geometry}
\usepackage{amsmath,amssymb,amsthm,mathtools,bm}
\usepackage{graphicx,microtype,booktabs,array,longtable,tabularx,multirow}
\usepackage{enumitem,caption,needspace,placeins}
\usepackage[table]{xcolor}
\usepackage[breakable]{tcolorbox}
\usepackage{tikz}
\usetikzlibrary{arrows.meta,positioning,calc}

% Default palette: the Okabe-Ito marks and light fills of them. Replace it to taste, in one place.
\definecolor{ink}{HTML}{222222}
\definecolor{structural}{HTML}{777777}
\definecolor{leader}{HTML}{B0B0B0}
\definecolor{markblue}{HTML}{0072B2}
\definecolor{markvermillion}{HTML}{D55E00}
\definecolor{markgreen}{HTML}{009E73}
\definecolor{markpurple}{HTML}{CC79A7}
\definecolor{markorange}{HTML}{E69F00}
\definecolor{marksky}{HTML}{56B4E9}
\definecolor{fillblue}{HTML}{D9EAF3}
\definecolor{fillvermillion}{HTML}{F9E7D9}
\definecolor{fillgreen}{HTML}{D9F0EA}
\definecolor{fillpurple}{HTML}{F7EBF2}
\definecolor{fillorange}{HTML}{FBF1D9}
\definecolor{fillsky}{HTML}{E6F4FC}

\usepackage[colorlinks=true,linkcolor=markblue,urlcolor=markgreen,citecolor=markgreen]{hyperref}
\usepackage[nameinlink,noabbrev]{cleveref}

% A procedure step, always in brackets: \procstep{Fit}.
\DeclareRobustCommand{\procstep}[1]{\mbox{\textup{[#1]}}}

% Thin body rules that only guide the eye. Top, header and bottom rules stay full booktabs.
\newcommand{\rowguide}{\arrayrulecolor{leader}\midrule[.15pt]\arrayrulecolor{ink}}
\newcommand{\rowguidepart}[1]{\arrayrulecolor{leader}\cmidrule[.15pt](lr){#1}\arrayrulecolor{ink}}

% Plain, fixed-size arrowheads, whose tips do not grow with line width.
\tikzset{>={Triangle[length=3.2pt,width=3pt]}}

% A recipe box holds only the short usable rule. Definitions and motivation go before it.
\newenvironment{recipe}[1]
  {\begin{tcolorbox}[breakable,colback=fillblue,colframe=markblue,boxrule=.5pt,arc=1pt,left=7pt,right=7pt,top=5pt,bottom=5pt,title={#1},fonttitle=\sffamily\bfseries]}
  {\end{tcolorbox}}

% Unfinished work: a design and a falsifiable expected outcome, never a fabricated result.
\newenvironment{resulttodo}[1]
  {\begin{tcolorbox}[breakable,colback=fillorange,colframe=structural,boxrule=.5pt,arc=1pt,left=7pt,right=7pt,top=5pt,bottom=5pt,title={Result to complete: #1},fonttitle=\sffamily\bfseries]}
  {\end{tcolorbox}}

\newtheoremstyle{statement}{\topsep}{\topsep}{\itshape}{}{\sffamily\bfseries}{.}{.5em}{}
\newtheoremstyle{exposition}{\topsep}{\topsep}{}{}{\sffamily\bfseries}{.}{.5em}{}
\theoremstyle{statement}
\newtheorem{theorem}{Theorem}[section]
\newtheorem{proposition}[theorem]{Proposition}
\newtheorem{lemma}[theorem]{Lemma}
\newtheorem{corollary}[theorem]{Corollary}
\theoremstyle{exposition}
\newtheorem{definition}[theorem]{Definition}
\newtheorem{conjecture}[theorem]{Conjecture}
\newtheorem{observation}[theorem]{Observation}
\newtheorem{example}[theorem]{Example}
\newtheorem{remark}[theorem]{Remark}

% An optional link from the tutorial's opening to the companion, and the anchors at the companion's statements.
\newcommand{\techfile}{{{SLUG}}-technical.pdf}
\newcommand{\techref}[2][companion]{\href{\techfile\#tech.#2}{#1}}
\newcommand{\techanchor}[1]{\hypertarget{tech.#1}{}}

\setlist{nosep,leftmargin=1.6em,topsep=4pt}
\setlength{\parindent}{1em}
\setlength{\parskip}{5pt plus 1pt minus 1pt}
\setlength{\emergencystretch}{1.5em}
\captionsetup{font=small,labelfont=bf,skip=5pt}
